\section{Directions that change slowly in scale and position}
\label{sec:labels}

The common prices in Proposition~\ref{prop:prices} pay for scalar
cancellation. Its remaining terms are nearby vector masses $a_q$.
We now choose directions so that these masses pay only a summable
amount for changes of direction. A maximal consecutive interval of levels
carrying the same direction will be called a \emph{run}. Ending a
run produces endpoint terms in the local estimates, so the first
cost is the probability of a change between successive levels.

There is also a spatial cost. We will average signed directional
estimates with weights depending on $x$ and then integrate by parts.
For any finite Lipschitz probability row $(\omega_j(x))_j$,
normalization gives
\[
 \sum_j\nabla\omega_j=0,
 \qquad
 \sum_j U_j\nabla\omega_j
   =\sum_j(U_j-U_*)\nabla\omega_j
\]
almost everywhere, for arbitrary scalar values $U_j$ and a common
value $U_*$. Thus the cost of differentiating a row is controlled
by the differences $U_j-U_*$, rather than the sizes of the values
themselves. In Section~\ref{sec:assembly}, the differences will be
bounded by sums of the small-radius oscillations $v_l$ defined in
Section~\ref{sec:prices}. Their integral over a level-$l$ cell costs
$r_l$ times its price and root masses, by \eqref{eq:prices}. A row
built from level-$k$ tents has derivatives of size $r_k^{-1}$.
The required joint bound therefore carries the ratio
$r_l/r_k=2^{k-l}$, for $k\le l$. This explains both the row-derivative
cost and its weight in the packing estimate below.

Retain $g\in C_c^\infty(\R^2)$, $h=\nabla g\,\dd x$,
$V=|h|(\R^2)$, and the levels $1,\ldots,m$, with $m\ge1$.
We use the hats, vector masses $h_q=a_q\theta_q$, and refinement
coefficients from Section~\ref{sec:prices}, defined at every integer
level. The coarser levels will be needed to compare nearby directions.
The palette $\mathcal N$ is the fixed $\tau/20$-net of the unit circle,
where $\tau=1/100$. Recall that $(q,n)$ is good when $\theta_q$
makes an angle at most $\pi/6$ with $n$.
All vertex and palette choices below use fixed tie rules, independently
of the observation point $x$. Directions at zero vector masses retain
the prescribed unit-vector convention.

\subsection{Directions in neighboring cells}

At each vertex we first look for a nearby direction carrying maximal
mass. This choice prevents a small neighboring mass from imposing an
unpaid change of direction. For $z\in V_k$, choose a vertex
$w(z)\in V_k$ maximizing $a_w$ among
$|w-z|_{\infty}\le300r_k$, and set
\[
 \Theta_z=\theta_{w(z)}.
\]
Choose $n_z\in\mathcal N$ with $|n_z-\Theta_z|\le\tau/20$.
The level is part of the vertex label in this notation. Define
\[
 E_k(q)=
 \max_{\substack{z\in V_k\\|z-q|_{\infty}\le200r_k}}
       |\Theta_z-\theta_q|^2,
 \qquad q\in V_k.
\]
In particular, $0\le E_k(q)\le4$. The quantity $E_k(q)$ measures
how far the nearby chosen directions can depart from its own
direction. Put $A_j=\sum_{p\in V_j}a_p$ for every integer $j$.
By \eqref{eq:prices-vector-mass-budget}, these total vector masses
satisfy $0\le A_j\le A_{j+1}\le V$. Their increase pays the
weighted sum of the disagreements.

\begin{lemma}[Neighboring direction budget]
\label{lem:labels-neighbors}
There is an absolute constant $C$ such that
\begin{equation}
 \sum_{k=1}^m\sum_{q\in V_k}a_q E_k(q)\le C V.
 \label{eq:neighbor-budget}
\end{equation}
\end{lemma}

\begin{proof}
Fix $k$ and $q$ with $a_q>0$, choose a maximizing $z$ in the definition
of $E_k(q)$, and put $w=w(z)$. Since $q$ is among the candidates for
$w(z)$, we have
\[
 a_w\ge a_q,
 \qquad |w-q|_{\infty}\le500r_k.
\]
Take $L=12$, and choose a nearest vertex $p\in V_{k-L}$ to $q$ in each
coordinate. Then
\[
 |q-p|_{\infty}\le\tfrac12r_{k-L},
 \qquad
 |w-p|_{\infty}
 \le\bigl(\tfrac12+500\,2^{-12}\bigr)r_{k-L}
 <\tfrac58r_{k-L}.
\]
Each coordinate factor of $\phi_p(q)$ is at least $1/2$, and each
coordinate factor of $\phi_p(w)$ is greater than $3/8$. Hence
$\lambda_{pq}\ge1/4$ and $\lambda_{pw}\ge9/64$.
For this coarse vertex write
\[
 H_{p,k}=\sum_{s\in V_k}\lambda_{ps}a_s-a_p.
\]
The vector refinement identity implies, also when $a_p=0$,
\begin{equation}
 H_{p,k}
 =\sum_{s\in V_k}\lambda_{ps}a_s
                  (1-\theta_p\cdot\theta_s)
 =\frac12\sum_{s\in V_k}\lambda_{ps}a_s
                  |\theta_s-\theta_p|^2.
 \label{eq:labels-neighbor-defect}
\end{equation}
Here $\theta_p\cdot h_p=a_p$ is true even at zero, for the prescribed
unit direction at zero. The triangle inequality for squared distances,
the two lower bounds on the interpolation coefficients, and $a_w\ge a_q$
therefore give
\[
 a_q E_k(q)
 \le2a_q\bigl(|\theta_q-\theta_p|^2
                    +|\theta_w-\theta_p|^2\bigr)
 \le C H_{p,k}.
\]
For fixed $k$, each coarse vertex is chosen by at most a fixed number
of fine vertices, depending only on $L$. Vertices with $a_q=0$ contribute
nothing. The column sums of the refinement coefficients now give
\[
 \sum_{q\in V_k}a_q E_k(q)
 \le C\sum_{p\in V_{k-L}}H_{p,k}
 =C(A_k-A_{k-L}).
\]
Finally,
\[
 \sum_{k=1}^m(A_k-A_{k-L})
 =\sum_{k=1}^m\sum_{j=k-L+1}^k(A_j-A_{j-1})
 \le L V,
\]
because every nonnegative increment appears at most $L$ times. This proves
\eqref{eq:neighbor-budget}.
\end{proof}

\subsection{The finite-state label law}

The neighboring budget does not prescribe a direction at every point.
We interpolate the fixed vertex choices by the tents, retaining the
previous direction unless it is appreciably far from the sampled one.
The interval on which the switching probability vanishes will keep
small direction fluctuations from creating repeated run endpoints.

For each $x\in\R^2$ we define a probability law $\Prob_x$ on
$\mathcal N^m$. At the first step, sample $z\in V_1$ with probabilities
$\phi_z(x)$ and set $n_1=n_z$. At a step $k\ge2$, condition on the old
label $n=n_{k-1}$, sample $z\in V_k$ with probabilities $\phi_z(x)$,
and switch to $n_z$ with probability
\begin{equation}
 p_k(n,z)=P(|n-\Theta_z|),
 \qquad
 P(d)=\min\{1,\tau^{-2}(d-\tau)_+^2\}.
 \label{eq:labels-switch-function}
\end{equation}
Otherwise retain $n$. These samples and switch trials are made independently
of the past conditional on the old label. If $p_k(n,z)>0$, then $n_z\ne n$:
the opposite equality would imply $|n-\Theta_z|\le\tau/20<\tau$.

For a label $n$, let $\delta_n$ denote the probability vector concentrated
at $n$. The initial row and subsequent transition rows are, respectively,
\begin{align}
 \pi_1(x)
 &=\sum_{z\in V_1}\phi_z(x)\delta_{n_z},
 \label{eq:labels-initial-row}\\
 \pi_k(x;n,\cdot)
 &=\delta_n+\sum_{z\in V_k}\phi_z(x)p_k(n,z)
                              (\delta_{n_z}-\delta_n),\quad k\ge2.
 \label{eq:labels-transition-row}
\end{align}
Every entry is nonnegative and every row sums to one, as is also clear
from the sampling rule. The rows are continuous, piecewise bilinear
interpolations of bounded vectors, with Lipschitz constant at most
$C/r_k$. These formulas, together with the initial row,
specify the probability of each label sequence as a product of entries;
this is the law $\Prob_x$. We denote its expectation by $\E_x$ and the
marginal law of $n_k$ by $\zeta_k(x)$. When convenient, the initial row is
viewed as a transition from a single auxiliary state $*$, with
$\zeta_0(x)=\delta_*$ for every $x$.

The row switch probability for a prescribed old label is
\[
 \rho_k(x\mid n)=\sum_{z\in V_k}\phi_z(x)p_k(n,z),\qquad k\ge2.
\]
Since a trial counted by this expression changes the label, define
\begin{equation}
 \begin{aligned}
 R_k(x)&=\sum_{n\in\mathcal N}\zeta_{k-1}(x;n)\rho_k(x\mid n)
       =\Prob_x(n_k\ne n_{k-1}),\quad 2\le k\le m,\\
 R_1(x)&=R_{m+1}(x)=1.
 \end{aligned}
 \label{eq:labels-switch-rates}
\end{equation}
The two sentinel values count the beginning and end of a history;
they apply also when $m=1$. The functions $R_k$ are continuous, since
the rows are Lipschitz and each marginal is obtained by finitely many
matrix products.

Let $G$ be the union of the grid lines at levels $1,\ldots,m$.
It is a fixed Lebesgue-null set. Off $G$ all the tent rows are
differentiable. Define the scaled expected row-derivative costs there by
\begin{equation}
 \begin{split}
 T_k(x)&=r_k\sum_{n\in\mathcal N}\zeta_{k-1}(x;n)
          \sum_{n'\in\mathcal N}
             |\nabla_x\pi_k(x;n,n')|,\quad 2\le k\le m,\\
 T_1(x)&=r_1\sum_{n'\in\mathcal N}
                         |\nabla_x\pi_1(x;n')|.
 \end{split}
 \label{eq:labels-derivative-costs}
\end{equation}
Set $T_k=0$ on $G$. These are measurable versions of the costs, and
all later bounds on their suprema are essential-supremum bounds.
The derivative in this definition acts on one transition row; it is
not the derivative of the marginal $\zeta_k$. Off $G$, at most four
level-$k$ tents have a nonzero gradient, and
$|\nabla\phi_z|\le\sqrt2/r_k$. It follows from the row formulas that
\begin{equation}
 0\le R_k\le1,
 \qquad 0\le T_k\le C.
 \label{eq:labels-uniform-costs}
\end{equation}
The norm inside the inner sum in \eqref{eq:labels-derivative-costs} is the
Euclidean norm of the spatial gradient, whereas norms of probability vectors
below are their $\ell^1$ norms.

\begin{lemma}[Good labels near a controlled vertex]
\label{lem:good-label}
Set $c_*=(\tau/10)^2$. Suppose $q\in V_k$, $E_k(q)\le c_*$, and
$|x-q|_{\infty}\le50r_k$. Then $(q,n_k)$ is good with $\Prob_x$-probability
one.
\end{lemma}

\begin{proof}
Every vertex $z$ with $\phi_z(x)>0$ satisfies
$|z-q|_{\infty}<51r_k<200r_k$, and hence
$|\Theta_z-\theta_q|\le\tau/10$. An initial or switched label therefore
has distance at most $\tau/10+\tau/20$ from $\theta_q$. If a label $n$ is
retained after sampling $z$ with positive probability, then
$P(|n-\Theta_z|)<1$, so $|n-\Theta_z|<2\tau$. Its distance from
$\theta_q$ is consequently less than $21\tau/10$.
Both bounds are less than $2\sin(\pi/12)$, since $\tau=1/100$.
The angle between the resulting label and $\theta_q$ is thus less than
$\pi/6$, as required by the definition of a good pair.
\end{proof}

This lemma pays the bad-root terms in the first line of
\eqref{eq:prices}. If $q\in\mathcal R(Q_p)$ at level $k$, then
$|x-q|_\infty\le21r_k$ throughout $Q_p$. Whenever $(q,n)$ is bad
and $E_k(q)\le c_*$, the marginal probability of $n_k=n$ therefore
vanishes on that whole cell. In the remaining case,
$a_q\le c_*^{-1}a_qE_k(q)$, whose sum is bounded by
\eqref{eq:neighbor-budget}, with the fixed cell and palette
multiplicities. It remains to pay root masses at the endpoints of
runs and in the terms where spatial weights are differentiated.

For $q\in V_l$ define the closed square
\[
 D(q)=\{x\in\R^2:|x-q|_\infty\le25r_l\}.
\]
It contains every level-$l$ pricing cell whose root neighborhood
contains $q$: indeed $q\in\mathcal R(Q_p)$ and $x\in Q_p$ give
$|x-q|_\infty\le21r_l$. We can now state the estimate that pays both
types of cost on those whole cells.

\begin{proposition}[Packing of switch and derivative costs]
\label{prop:label-packing}
For every real $g\in C_c^\infty(\R^2)$ and every $m\ge1$, the
quantities constructed above satisfy
\begin{equation}
 \sum_{l=1}^m\sum_{q\in V_l}a_q
 \left[
   \sup_{x\in D(q)}\bigl(R_l(x)+R_{l+1}(x)\bigr)
   +\sum_{k=1}^l2^{k-l}
       \operatorname*{ess\,sup}_{x\in D(q)}T_k(x)
 \right]
 \le C V.
 \label{eq:label-packing}
\end{equation}
The constant is absolute and independent of $m$.
\end{proposition}

We prove this proposition in three steps. First a positive measure
on refinement paths will realize each $a_q$ exactly. Along each path,
vector additivity will bound switches at its terminal position.
Finally comparison of the laws at nearby points will turn that
fixed-position estimate into the suprema over $D(q)$.

\subsection{Particles and intrinsic direction changes}

The scalar decomposition in Section~\ref{sec:prices} is an identity
of signed measures with cancelling pairs. Here we need a different
object: a positive measure on complete refinement paths, with exact
occupancy $a_q$ after birth. It must also retain the path before birth,
because the label law has been evolving since level one.
This path measure will not depend on $x$. We will combine it with
$\Prob_x$ by evaluating that law at a path's terminal vertex.
A path node at level $k$ is
$\xi=(q_1,\ldots,q_k)$, where $q_j\in V_j$ and each consecutive
refinement coefficient is positive. Its path weight and mass are
\[
 b_\xi=\prod_{j=1}^{k-1}\lambda_{q_jq_{j+1}},
 \qquad M_\xi=b_\xi a_{q_k}.
\]
The empty product at a root is one. The associated vector is
$b_\xi h_{q_k}$. Vectors add from the children to their parent, and therefore
\begin{equation}
 M_\xi\le\sum_{\eta\text{ child of }\xi}M_\eta.
 \label{eq:labels-child-capacity}
\end{equation}
Only a finite forest is needed. Choose all level-one roots whose closed
tent supports meet $\supp h$, and retain all their positive-weight paths
through level $m$. There are finitely many roots and each node has finitely
many children. If $a_q>0$, the support of $\phi_q$ meets $\supp h$.
Every ancestor tent then meets $\supp h$ by support nesting, so every
positive-weight path to $q$ starts among the chosen roots. In
particular no contribution to a nonzero mass is lost by this finite
restriction. If $h=0$, the forest may be empty.
Consequently, at each level $k$,
\begin{equation}
 \sum_{\xi:\,q_k=q}b_\xi=1\quad\text{if }a_q>0,
 \qquad
 \sum_{|\xi|=k}M_\xi=A_k.
 \label{eq:labels-node-column-sums}
\end{equation}

\begin{lemma}[A finite particle measure]
\label{lem:labels-particles}
There is a finite positive measure $\mu$ on pairs consisting of a terminal
path $(q_1,\ldots,q_m)$ and a birth level $B\in\{1,\ldots,m\}$ such that,
for every node $\xi$ at level $k$,
\begin{equation}
 \mu\{B\le k,\ (q_1,\ldots,q_k)=\xi\}=M_\xi.
 \label{eq:labels-node-occupancy}
\end{equation}
In particular,
\begin{equation}
 \mu\{B\le k,\ q_k=q\}=a_q,
 \qquad \mu(\text{all particles})=A_m\le V.
 \label{eq:labels-vertex-occupancy}
\end{equation}
Writing $y=q_m$ for the terminal position of a particle, its complete path
obeys
\begin{equation}
 |y-q_s|_{\infty}<r_s\qquad(1\le s\le m).
 \label{eq:labels-path-locality}
\end{equation}
The path is defined also at levels preceding its birth.
\end{lemma}

\begin{proof}
At each root create mass $M_\xi$ with birth level one. Suppose the population
already present at a level-$k$ node has mass $M_\xi$, and put
$S_\xi=\sum_{\eta\text{ child}}M_\eta$. If $S_\xi>0$, send the fraction
$M_\eta/S_\xi$ of each existing population type to the child $\eta$.
The mass sent there is $M_\xi M_\eta/S_\xi\le M_\eta$ by
\eqref{eq:labels-child-capacity}. Create at that child additional mass
$M_\eta-M_\xi M_\eta/S_\xi$ with birth level $k+1$. If $S_\xi=0$, all
these masses are zero. This operation preserves existing particles and
fills each child to exactly its prescribed mass. A new particle at a
child ending in $q_{k+1}$ receives the entire
prefix $(q_1,\ldots,q_k,q_{k+1})$ of that child, as well as its new
birth level. Thus a particle born below a zero-mass ancestor still
has a specified direction history at that ancestor. No particle
mass is counted as occupying that ancestor before birth.

After the finite number of levels, the masses of the resulting
terminal-path and birth-level types define a finite atomic measure
$\mu$. All particles present at $\xi$ continue to terminal descendants
of $\xi$, conserving their mass $M_\xi$. Conversely, any terminal
particle with that prefix and $B\le k$ was already present at $\xi$.
Particles born later may carry the same prefix, but the condition
$B\le k$ excludes them. This proves \eqref{eq:labels-node-occupancy}.
Summing over prefixes and using \eqref{eq:labels-node-column-sums} gives
\eqref{eq:labels-vertex-occupancy}.
For a positive consecutive refinement coefficient,
$|q_{j+1}-q_j|_{\infty}\le r_j/2$. Hence
\[
 |q_m-q_s|_{\infty}
 \le\sum_{j=s}^{m-1}r_j/2
 =r_s(1-2^{-(m-s)})<r_s,
\]
with the same strict upper bound when $s=m$. This proves the locality
assertion.
\end{proof}

The particle measure now realizes every vector mass exactly at its
level. It lets us prove a pathwise bound and then recover a sum over
all cells by integration. We next mark the genuine direction changes
along a path. The sum of the norms of descendant vectors may exceed
the norm of their sum at the ancestor; this nonnegative difference
will pay for the marks.

Mark every root of the path forest. Below a mark, mark the first nodes
whose direction $\theta_{q_s}$ has distance greater than $\tau/8$ from
the direction at that mark, and repeat the rule below each new mark.
These intrinsic marks depend on the path forest and the vector masses;
they do not depend on the label process at any point $x$.

\begin{lemma}[Intrinsic marks and expected switches]
\label{lem:labels-switch-packing}
The particle measure satisfies
\begin{align}
 \sum_{\xi\text{ intrinsically marked}}M_\xi&\le C V,
 \label{eq:labels-mark-budget}\\
 \int\sum_{s=B}^m E_s(q_s)\,\dd\mu&\le C V,
 \label{eq:labels-particle-energy}\\
 \int\sum_{s=B}^m R_s(y)\,\dd\mu&\le C V.
 \label{eq:labels-particle-switches}
\end{align}
\end{lemma}

\begin{proof}
For a marked nonterminal node $\xi$, let $\mathcal L_\xi$ consist of
the first subsequent marks on each descending branch, supplemented by
terminal nodes on branches with no later mark. This is a finite frontier.
Writing $\vartheta_\xi=\theta_{q_k}$ at a level-$k$ node, vector additivity
gives
\begin{equation}
 \sum_{\eta\in\mathcal L_\xi}M_\eta-M_\xi
 =\frac12\sum_{\eta\in\mathcal L_\xi}M_\eta
                   |\vartheta_\eta-\vartheta_\xi|^2
 \ge\frac{\tau^2}{128}
          \sum_{\substack{\eta\in\mathcal L_\xi\\\eta\text{ marked}}}
                   M_\eta.
 \label{eq:labels-stopped-defect}
\end{equation}
For clarity, vector additivity on the frontier says
$\sum_{\eta\in\mathcal L_\xi}M_\eta\vartheta_\eta
=M_\xi\vartheta_\xi$. Taking its scalar product with the unit vector
$\vartheta_\xi$ proves the identity, even if $M_\xi=0$.
For $m>1$, every marked nonterminal node other than a root occurs positively
in its preceding marked frontier and negatively in its own
summand. Every terminal node occurs in exactly one frontier.
The sum of the left sides is therefore $A_m-A_1\le V$. Every nonroot mark
occurs in exactly one frontier. Adding the root masses proves
\eqref{eq:labels-mark-budget}; if $m=1$, the assertion follows directly
from the root mass bound.

By \eqref{eq:labels-node-occupancy}, integration of the number of intrinsic
marks at levels $s\ge B$ is the left side of
\eqref{eq:labels-mark-budget}. Similarly,
\[
 \int\sum_{s=B}^m E_s(q_s)\,\dd\mu
 =\sum_{s=1}^m\sum_{q\in V_s}a_q E_s(q),
\]
so \eqref{eq:labels-particle-energy} follows from
\eqref{eq:neighbor-budget}.

For the switch estimate, fix a particle, including its path, birth level,
and terminal position $y$. Run the finite-state law at this fixed point
$y$. Let $X_1=1$, and for $s\ge2$ let
$X_s=\1_{\{n_s\ne n_{s-1}\}}$. Thus $\E_y X_s=R_s(y)$ for $1\le s\le m$.
Among levels $B,\ldots,m$, call a level exceptional if it is intrinsically
marked or if $E_s(q_s)>c_*$. Remove those levels and consider a consecutive
interval of remaining levels. Throughout this interval, all
$\theta_{q_s}$ are within $\tau/8$ of the same most recent intrinsic
marked direction, even if that mark occurred before birth. Every sampled
vertex at $y$ lies within $2r_s$ of $q_s$ by
\eqref{eq:labels-path-locality}, and hence its sampled direction is within
$\tau/10$ of $\theta_{q_s}$. Any two sampled directions on the interval
are therefore at distance at most
\[
 2(\tau/8+\tau/10)=9\tau/20.
\]
After one switch on the interval, the new label is within an additional
$\tau/20$ of the direction sampled at that switch. Its distance to every
subsequently sampled direction is at most $\tau/2<\tau$, so all later
switch probabilities on this interval are zero. The same observation
applies after the initial assignment, should it be counted on such an
interval. There is consequently at most one counted switch on each
remaining interval, regardless of the label entering it.

If there are $J$ exceptional levels, there are at most $J+1$ such
intervals and at most $J$ switches at exceptional levels. Thus every
realization obeys
\[
 \sum_{s=B}^m X_s
 \le1+2\#\{s\ge B:s\text{ is intrinsically marked}\}
       +2c_*^{-1}\sum_{s=B}^m E_s(q_s).
\]
Taking $\E_y$, then integrating over particles, proves
\eqref{eq:labels-particle-switches} from the first two bounds and
$\mu(\text{all particles})\le V$.
\end{proof}

\subsection{Comparison of label laws at nearby points}

We have bounded the expected switches at the single point $y$ attached
to each particle. This does not yet give Proposition~\ref{prop:label-packing}:
the local prices are integrated on cells, so their coefficients must
be controlled uniformly there. We also still need the row-derivative
cost. The same comparison of nearby rows supplies both missing bounds.

We compare the law at any point of $D(q_l)$ with the law at the
terminal position of a particle. The comparison uses the complete particle
path, including its levels before birth.

\begin{lemma}[Spatial comparison]
\label{lem:labels-spatial-comparison}
Fix a particle, a level $1\le l\le m$, and $x\in D(q_l)$. Set
\[
 c_s=R_s(y)+E_s(q_s),\qquad 1\le s\le m.
\]
For $1\le k\le\min(l+1,m)$,
\begin{equation}
 R_k(x)+T_k(x)
 \le C\sum_{s=1}^k2^{s-k}c_s.
 \label{eq:labels-spatial-comparison}
\end{equation}
The assertion for $R_k$ holds at every $x$ in the specified set, and the
assertion for $T_k$ holds almost everywhere there. The constant is uniform
over particles, $l$, and $x$.
\end{lemma}

\begin{proof}
By \eqref{eq:labels-path-locality},
$|x-y|_{\infty}<26r_l$. At a step $s\le\min(l+1,m)$, any vertex $z$ whose
tent is nonzero at $x$ or at $y$ satisfies
\[
 |z-q_s|_{\infty}
 \le r_s+26r_l+r_s
 \le54r_s<200r_s.
\]
Off $G$, a tent with nonzero gradient at $x$ also has positive value
there, so the same bound covers every differentiated tent. At grid
points the probability comparison still uses the exact tent values;
no assertion about a classical row derivative is needed.
Thus all the directions involved differ from
$\theta_{q_s}$ by at most $\sqrt{E_s(q_s)}$.

Fix $s\ge2$ and an old label $n$. The function
$d\mapsto\sqrt{P(d)}=\min\{1,\tau^{-1}(d-\tau)_+\}$ is
$\tau^{-1}$-Lipschitz. For any two involved vertices $z,z'$, it follows that
\[
 p_s(n,z)
 \le2p_s(n,z')+2\tau^{-2}|\Theta_z-\Theta_{z'}|^2
 \le2p_s(n,z')+8\tau^{-2}E_s(q_s).
\]
Averaging $z'$ with the probabilities $\phi_{z'}(y)$ shows that the largest
$p_s(n,z)$ among all involved vertices is bounded by
\begin{equation}
 C\bigl(\rho_s(y\mid n)+E_s(q_s)\bigr).
 \label{eq:labels-row-rate-comparison}
\end{equation}
The row formula \eqref{eq:labels-transition-row} and the tent derivative
bounds then give
\begin{align}
 \rho_s(x\mid n)
 &\le C\bigl(\rho_s(y\mid n)+E_s(q_s)\bigr),
 \label{eq:labels-row-rate}\\
 r_s\sum_{n'\in\mathcal N}|\nabla_x\pi_s(x;n,n')|
 &\le C\bigl(\rho_s(y\mid n)+E_s(q_s)\bigr)
 \label{eq:labels-row-derivative}
\end{align}
almost everywhere in the second assertion. Moreover, the union of the
sets of tents nonzero at $x$ and $y$ has uniformly bounded cardinality,
and each tent is $C/r_s$-Lipschitz. Subtracting the two rows therefore
yields the pointwise bound
\begin{equation}
 \|\pi_s(x;n,\cdot)-\pi_s(y;n,\cdot)\|_{\ell^1}
 \le C2^{s-l}\bigl(\rho_s(y\mid n)+E_s(q_s)\bigr).
 \label{eq:labels-row-difference}
\end{equation}
For $s=1$, the analogous derivative and difference bounds are $C$ and
$C2^{1-l}$, respectively, by \eqref{eq:labels-initial-row}.
These bounds are covered by using $c_1=1+E_1(q_1)\ge1$.

The rows have now been compared for each fixed old label. The old-label
distributions themselves also differ between $x$ and $y$, and must be
accounted for before estimating $R_k$ and $T_k$.
Let $d_s(x,y)=\|\zeta_s(x)-\zeta_s(y)\|_{\ell^1}$ and $d_0=0$.
Multiplication by a stochastic matrix $\Pi$ contracts the $\ell^1$
norm of signed row vectors: for any such vector $v$,
\[
 \|v\Pi\|_1
 \le\sum_n |v(n)|\sum_{n'}\Pi(n,n')=\|v\|_1.
\]
The marginal recursion therefore gives
\begin{align*}
 d_s(x,y)
 &\le d_{s-1}(x,y)
   +\|\zeta_{s-1}(y)(\pi_s(x)-\pi_s(y))\|_{\ell^1}\\
 &\le d_{s-1}(x,y)+C2^{s-l}c_s.
\end{align*}
The first step uses the single-row interpretation. At later steps the
last inequality follows by averaging \eqref{eq:labels-row-difference}
against the old-label law at $y$, which averages
$\rho_s(y\mid n)$ to $R_s(y)$. Consequently
\begin{equation}
 d_j(x,y)\le C\sum_{s=1}^j2^{s-l}c_s,
 \qquad j\le\min(l+1,m).
 \label{eq:labels-marginal-distance}
\end{equation}
This recursion does not assume independence between marginal labels
at the two observation points.

For $k\ge2$, first average the row switch rate and the row derivative
cost at $x$ against $\zeta_{k-1}(y)$. Equations
\eqref{eq:labels-row-rate}--\eqref{eq:labels-row-derivative} bound the
result by $C c_k$. Replacing this old-label distribution by
$\zeta_{k-1}(x)$ costs at most $C d_{k-1}(x,y)$, because row switch rates
are at most one and scaled row derivative costs are uniformly bounded.
Thus
\[
 R_k(x)+T_k(x)
 \le C c_k+C\sum_{s=1}^{k-1}2^{s-l}c_s.
\]
Since $k\le l+1$, we have $2^{s-l}\le2\,2^{s-k}$. This proves
\eqref{eq:labels-spatial-comparison} for $k\ge2$. For $k=1$ it follows
directly from $R_1=1$, \eqref{eq:labels-uniform-costs}, and $c_1\ge1$.
\end{proof}

\subsection{The packing estimate}

\begin{proof}[Proof of Proposition~\ref{prop:label-packing}]
Denote the bracket for a level-$l$ vertex by $\mathcal C_l(q)$.
All its quantities are nonnegative and measurable; the switch rates are
continuous and the derivative costs have specified measurable versions.
The particle occupancy identity gives the exact rewriting
\begin{equation}
 \sum_{l=1}^m\sum_{q\in V_l}a_q\mathcal C_l(q)
 =\int\sum_{l=B}^m\mathcal C_l(q_l)\,\dd\mu.
 \label{eq:labels-cost-occupancy}
\end{equation}
Fix a particle, and continue to write $c_s=R_s(y)+E_s(q_s)$.
The comparison in Lemma~\ref{lem:labels-spatial-comparison} is uniform on
each $D(q_l)$, so it applies to the supremum and essential supremum in
$\mathcal C_l(q_l)$. For $l<m$, it bounds both $R_l$ and $R_{l+1}$.
For $l=m$, the convention $R_{m+1}=1$ contributes just one additional
unit. The derivative terms have the exact geometric convolution
\[
 \sum_{k=1}^l2^{k-l}\sum_{s=1}^k2^{s-k}c_s
 =\sum_{s=1}^l(l-s+1)2^{s-l}c_s.
\]
It follows that
\begin{equation}
 \sum_{l=B}^m\mathcal C_l(q_l)
 \le C\left[
  1+\sum_{l=B}^m
     \sum_{s=1}^{\min(l+1,m)}(l-s+2)2^{s-l}c_s
 \right].
 \label{eq:labels-double-convolution}
\end{equation}

For a fixed $s\ge B$, the coefficient of $c_s$ on the right is bounded by
\[
 \sum_{j=-1}^{\infty}(j+2)2^{-j}<\infty,
 \qquad j=l-s.
\]
Thus the terms at or after birth cost at most
$C\sum_{s=B}^m c_s$. Their integral is at most $C V$ by
\eqref{eq:labels-particle-energy} and
\eqref{eq:labels-particle-switches}.

For the terms before birth, use $0\le c_s\le5$. Put $q=B-s\ge1$ and
$j=l-B\ge0$. Their total for this particle is at most
\[
 5\sum_{q=1}^{\infty}\sum_{j=0}^{\infty}
           (j+q+2)2^{-(j+q)}
 =5\sum_{n=1}^{\infty}n(n+2)2^{-n}<\infty.
\]
This bound is independent of $B$ and of the particle, and therefore its
integral is at most $C V$ by \eqref{eq:labels-vertex-occupancy}.
The isolated unit in \eqref{eq:labels-double-convolution} has the same
bound. Inserting these estimates into \eqref{eq:labels-cost-occupancy}
proves \eqref{eq:label-packing}.
\end{proof}
